\section{A finite Wronskian formula for every Hermite orbital radius}
\label{sec:analytic-orbital-wronskians}

The constants in NN-A029 admit a one-variable description. The
critical-length theorem used below is a classical theorem of extended
Chebyshev spaces; the identification with the Hermite orbital radius
and the resulting detecting constants are the deductions here.

For $m\ge2$ put
\begin{equation}\label{eq:analytic-hermite-ode}
 P_m(z)=i^m\operatorname{He}_m(z/i)
 =\sum_{j=0}^{\lfloor m/2\rfloor}
       \frac{m!}{2^j j!(m-2j)!}\,z^{m-2j}.
\end{equation}
Its distinct roots are $\alpha_j=ih_{m,j}$. Define the real entire
function
\begin{equation}\label{eq:analytic-fundamental-solution}
 S_m(x)=\sum_{j=1}^m\frac{e^{\alpha_jx}}{P_m'(\alpha_j)}
\end{equation}
and the Hankel Wronskians
\begin{equation}\label{eq:analytic-finite-wronskians}
 W_{m,k}(x)=\det\bigl[S_m^{(p+q)}(x)\bigr]_{p,q=0}^k,
 \qquad 0\le k\le\left\lfloor\frac{m-2}{2}\right\rfloor.
\end{equation}
The first positive zero of a function with no positive zero is, by
convention, $+\infty$.

\begin{theorem}[NN-A030: exact Wronskian reduction and computability]
\label{thm:analytic-orbital-wronskians}
The orbital radius of NN-A029 is
\begin{equation}\label{eq:analytic-wronskian-radius}
 \boxed{\displaystyle
 \mathfrak r_m=
 \min_{0\le k\le\lfloor(m-2)/2\rfloor}
        \inf\{x>0:W_{m,k}(x)=0\}.}
\end{equation}
Equivalently, it is the critical interval length for Hermite
interpolation in the real solution space of $P_m(D)f=0$.
In particular,
\begin{equation}\label{eq:analytic-radius-elementary-bounds}
       \frac\pi{h_m}\le\mathfrak r_m
                         \le\frac{(m-1)\pi}{h_m}.
\end{equation}
There is an algorithm, uniform in $m$, which returns rational bounds
$c<\mathfrak r_m<C$ of any prescribed positive rational width.
This is a computability assertion, without a polynomial-time or
uniform practical running-time bound. Consequently every local
constant in NN-A029 is specified by finitely many explicit
one-variable entire functions. Determining the actual endpoint
multiplicities remains a separate problem.
\end{theorem}

\begin{proof}
\emph{The differential equation and its canonical basis.}
The coefficient formula in \eqref{eq:analytic-hermite-ode} follows
directly from the defining formula for the probabilists' Hermite
polynomial. Thus $P_m$ is a real monic polynomial of parity $m$.
Complex conjugation permutes the terms in
\eqref{eq:analytic-fundamental-solution}, so $S_m$ is real on the
real axis. The elementary Lagrange interpolation identities give
\[
 \sum_j\frac{\alpha_j^r}{P_m'(\alpha_j)}=0
       \quad(0\le r<m-1),\qquad
 \sum_j\frac{\alpha_j^{m-1}}{P_m'(\alpha_j)}=1.
\]
Hence $S_m$ is the unique solution of $P_m(D)S_m=0$ with
\[
 S_m^{(r)}(0)=0\quad(0\le r<m-1),\qquad
 S_m^{(m-1)}(0)=1.
\]
Let $E_m$ be the real solution space. It has dimension $m$, and is
closed under translations and reflections. The functions
\[
 u_j=S_m^{(m-1-j)},\qquad 0\le j\le m-1,
\]
form a canonical basis: $u_j^{(r)}(0)=0$ for $r<j$ and
$u_j^{(j)}(0)=1$. Their complex span equals the span of
$e^{\alpha_1x},\ldots,e^{\alpha_mx}$.

\emph{Orbital determinants and arbitrary Hermite nodes.}
Write $Q=m(m-1)/2$ and $C_m=\prod_{q=1}^{m-1}q!$. The normalized
HCIZ identity, with both Vandermondes in the same order, says
\begin{equation}\label{eq:analytic-confluent-hciz}
 \phi_m(\operatorname{diag}(x_1,\ldots,x_m))
 =C_m\frac{\det[e^{\alpha_jx_k}]_{j,k=1}^m}
               {\Delta(\alpha)\Delta(x)},
 \qquad \Delta(\alpha)=i^Q\Delta(h).
\end{equation}
The identity extends continuously to repeated $x_k$. For distinct
nodes $t_1<\cdots<t_s$ with multiplicities $\nu_1+\cdots+\nu_s=m$,
the confluent numerator has columns
\[
 \left(\frac{\alpha_j^r e^{\alpha_jt_\ell}}{r!}\right)_{j=1}^m,
 \qquad 1\le\ell\le s,\quad0\le r<\nu_\ell,
\]
and the remaining denominator is the nonzero constant
$\Delta(\alpha)\prod_{\ell<v}(t_v-t_\ell)^{\nu_\ell\nu_v}$.
Thus the orbital value vanishes exactly when this Hermite
evaluation matrix is singular. Changing from the complex
exponential basis to a real basis of $E_m$ shows that singularity
is equivalent to a nonzero real $f\in E_m$ with zeros of the
prescribed multiplicities. In particular, this includes all
nonreduced node configurations, rather than just distinct nodes.

An $m$-dimensional real space is an extended Chebyshev space on an
interval if every nonzero member has at most $m-1$ zeros there,
counting multiplicities. By NN-A029, $\phi_m$ is strictly positive
for every spectrum of diameter less than $\mathfrak r_m$.
Consequently $E_m$ is an extended Chebyshev space on every
$[0,c]$ with $c<\mathfrak r_m$.

Conversely, if $c>\mathfrak r_m$, choose a negative orbital value
with spectral diameter less than $c$. Subtracting a scalar matrix
does not change it, since $\operatorname{Tr}D_m=0$. Along the
radial path from the zero matrix to this spectrum, the orbital
value starts at $1$ and ends negative. It vanishes at an
intermediate spectrum lying in $[0,c]$. The confluent determinant
criterion therefore makes $E_m$ fail the extended Chebyshev
property there. At the endpoint $c=\mathfrak r_m$, take negative
spectra with diameters decreasing to $\mathfrak r_m$, normalize
their minimum to zero, and extract a convergent subsequence. The
limit has orbital value zero: a negative limit could be radially
contracted to contradict the infimum. The endpoint also fails.
This proves that the critical interval length is exactly
$\mathfrak r_m$.

\emph{The finite Wronskian theorem and its precise indices.}
Carnicer--Mainar--Pe\~na, \emph{Spherical Bessel functions and
critical lengths}, RACSAM \textbf{117} (2023), article~29,
section~3, Proposition~3, p.~6, states that for a
reflection-invariant finite-dimensional space with a canonical
basis $(u_0,\ldots,u_n)$ and finite critical length, that length is
the least positive zero among
\[
                 \det W(u_j,\ldots,u_n),\qquad j>n/2.
\]
All hypotheses hold here; finiteness was proved in NN-A029 and
also follows from the elementary upper bound below. Put $n=m-1$
and $k=m-1-j$. Reversing these Wronskian columns changes only a
constant sign and gives $W_{m,k}$. The condition $j>(m-1)/2$
becomes $0\le k\le\lfloor(m-2)/2\rfloor$. This proves
\eqref{eq:analytic-wronskian-radius}. There are exactly
$\lfloor m/2\rfloor$ functions to test.

The lower bound in \eqref{eq:analytic-radius-elementary-bounds}
is NN-A029. For the upper bound, the nonzero function
$\sin(h_mx)$ belongs to $E_m$ and has the $m$ distinct zeros
\[
                0,\ \pi/h_m,\ \ldots,\ (m-1)\pi/h_m.
\]
The extended Chebyshev property therefore fails on their span.

\emph{Effective two-sided approximation, including tangential
Wronskian zeros.}
All roots of $\operatorname{He}_m$ are distinct real algebraic
numbers; isolating them and evaluating the finite exponential
sums with rigorous errors are effective. At zero,
\begin{equation}\label{eq:analytic-wronskian-leading}
 W_{m,k}(x)=(-1)^{k(k+1)/2}
 \left(\prod_{j=0}^k\frac{j!}{(m-1-j)!}\right)
 x^{(k+1)(m-1-k)}(1+O(x)).
\end{equation}
This follows by inserting the leading monomials
$S_m^{(j)}(x)=x^{m-1-j}/(m-1-j)!+O(x^{m-j})$ and evaluating
their Vandermonde Wronskian. Divide by the displayed power of $x$
and its nonzero signed leading coefficient. The resulting entire
functions have value $1$ at zero, computable Taylor coefficients
and effective uniform approximation bounds on every bounded
interval. Here are explicit bounds. Set $q=k+1$,
$d=q(m-1-k)$, and let $c_{m,k}$ denote the signed nonzero leading
coefficient in \eqref{eq:analytic-wronskian-leading}.
Cauchy--Binet gives the finite exponential sum
\[
 W_{m,k}(x)=\sum_{|J|=q} A_J e^{b_Jx},\qquad
 A_J=\frac{\Delta(\alpha_J)^2}{\prod_{j\in J}P_m'(\alpha_j)},
 \quad b_J=\sum_{j\in J}\alpha_j.
\]
Put $B=\sum_J|A_J|$ and $M=qh_m$. All $b_J$ are purely imaginary
and $|b_J|\le M$. For $T(x)=W_{m,k}(x)/(c_{m,k}x^d)$, its degree-$N$
Taylor polynomial $T_N$ therefore satisfies, on $|x|\le R$,
\[
 |T-T_N|\le\frac{BM^d}{|c_{m,k}|}
      \frac{(MR)^{N+1}e^{MR}}{(N+d+1)!},\qquad
 \sup_{|x|\le R}|T'(x)|
             \le\frac{BM^{d+1}}{|c_{m,k}|}e^{MR}.
\]
Indeed the coefficients of $T$ are
$W_{m,k}^{(n+d)}(0)/(c_{m,k}(n+d)!)$ and their absolute values are at
most $BM^{n+d}/(|c_{m,k}|(n+d)!)$. Summing this majorant gives both
bounds. The algebraic numbers $A_J,b_J$ and the elementary
exponentials allow outward rational bounds for every displayed
quantity.

For any rational $c<\mathfrak r_m$, all these finitely many
normalized Wronskians are strictly positive on $[0,c]$. Their
positive minima can therefore be certified by successive
uniform approximations and a finite mesh with rigorous
derivative bounds. This provides a semidecision procedure for
valid rational lower bounds. A certificate cannot pass at or
beyond the first zero.

Independently enumerate distinct rational spectra and certify
negative orbital values using \eqref{eq:analytic-confluent-hciz}.
Every rational $C>\mathfrak r_m$ admits such a negative spectrum
of diameter less than $C$: a negative spectrum exists by the
definition of the infimum, and negativity persists under a
sufficiently small perturbation to distinct rational entries.
This supplies certified upper bounds. Dovetail the lower and
upper searches until their width is smaller than the requested
error. Both searches terminate for bounds chosen sufficiently
close to the radius. No decision about whether an individual
Wronskian root is multiple, and no unjustified sign-change test
at that root, is used. This proves the algorithmic claim.
\end{proof}

\paragraph{The first examples.}
Here $S_2(x)=\sin x$ and
$S_3(x)=(1-\cos(\sqrt3x))/3$, so the single-Wronskian cases give
$\mathfrak r_2=\pi$ and $\mathfrak r_3=2\pi/\sqrt3$ again.
For $m=4$, put
\[
 a=\sqrt{3-\sqrt6},\qquad b=\sqrt{3+\sqrt6},\qquad
 S_4(x)=\frac{\sin(ax)/a-\sin(bx)/b}{b^2-a^2}.
\]
Then $\mathfrak r_4$ is exactly the smaller of the first positive
zeros of $S_4$ and $S_4S_4''-(S_4')^2$. This is a finite scalar
root problem, without a numerical value being asserted here.

\paragraph{A checked source correction.}
Beccari--Casciola--Mazure, \emph{Critical length: an alternative
approach}, arXiv:1904.09253v1, p.~12, Theorem~4.3 and
Corollary~4.4, equations~(19)--(20), print $W\ne0$ inside the
infima. The original page image was checked: this is a printed
sign error, not an extraction artifact. Those conditions must
read $W=0$, as the preceding sentence's first-zero description
and the correct Proposition~3 above specify. For $P_2(D)=D^2+1$
the printed nonzero condition gives infimum $0$ from $S_2=\sin x$,
whereas the critical length is $\pi$. No claim is made here
about whether later versions corrected that display.

\paragraph{Geometric translation.}
The algebra $\mathbb Q[z]/(P_m)$ is finite etale because $P_m$
has distinct roots; the finite-free discriminant criterion is
SGA~1, Expos\'e~I, section~4, Proposition~10. A companion matrix
$B_m$ of $P_m$ gives the actual algebraic connection $d-B_m\,dx$ on a
rank-$m$ trivial bundle on $\mathbb A^1$. Its analytic horizontal
sections are the solution space above. Evaluation on the finite
nonreduced scheme
\[
              \coprod_{\ell=1}^s
                  \operatorname{Spec}\mathbb C[x]/(x-t_\ell)^{\nu_\ell}
\]
is exactly the Hermite jet map, and its determinant vanishes
precisely when interpolation fails. A zero of $W_{m,k}(x)$ gives
a nonzero solution vanishing to order at least $m-1-k$ at $0$
and at least $k+1$ at $x$: use the span of
$u_{m-1-k},\ldots,u_{m-1}$. Thus the Wronskian formula reduces
this real interval boundary to finitely many two-point jet tests.
It does not turn the exponential connection into an arithmetic
Frobenius object, a pure motive, or a polynomial-time algorithm.
